% ======================================================
%  PART III — TỔNG QUAN GIẢI PHÁP
% ======================================================
% Yêu cầu: Mô tả tổng quan về concept giải pháp.
% - Form factor: smartphone app? tablet app? web page? smartwatch?
% - Interaction model: Users sẽ sử dụng/thao tác như thế nào?
% - Giá trị cốt lõi: Giải pháp này giải quyết pain points ra sao?
\section{Overall Solution Concept}

% --- Hướng dẫn viết ---
% Phần này KHÔNG đi vào chi tiết — chỉ tổng quan.
% Chi tiết từng conceptual solution đã trình bày ở Part II.

\subsection{Nền tảng và hình thức}

Giải pháp của nhóm là một trang web phát trực tuyến nội dung thể thao, được thiết kế để sử dụng trên máy tính để bàn hoặc laptop. Giao diện tận dụng lợi thế về kích thước màn hình lớn của máy tính để hiển thị đồng thời nhiều thông tin mà không gây quá tải cho người dùng -- điều khó thực hiện trên điện thoại khi diện tích màn hình bị hạn chế.

Trang web được chọn làm nền tảng chính vì một số lý do: Thứ nhất, đây là lựa chọn ``best of both worlds'' -- vừa có không gian màn hình rộng hơn điện thoại, vừa dễ thao tác hơn tivi (người dùng có chuột và bàn phím để điều hướng nội dung một cách linh hoạt). Thứ hai, trang web không yêu cầu người dùng cài đặt ứng dụng riêng, giúp giảm rào cản tiếp cận với người dùng.

\subsection{Cách người dùng tương tác với hệ thống}

Để giải quyết vấn đề trì hoãn khi xem trực tiếp, hệ thống\footnote{Chi tiết được trình bày ở Phần II, Vấn đề 1.} sử dụng trí tuệ nhân tạo để đánh dấu những khoảnh khắc quan trọng trong trận đấu. Người dùng có thể nhảy đến ngay các khoảnh khắc này để xem lại nếu bỏ lỡ. Bên cạnh đó, tính năng \textit{chế độ nhanh} cho phép hiển thị hình ảnh tĩnh kèm văn bản mô tả cho các đoạn không quan trọng, giúp tiết kiệm dữ liệu. Về lâu dài, mạng lưới P2P có thể được tích hợp để giảm tải cho máy chủ trong giờ cao điểm.

Để giải quyết vấn đề quá tải thông tin, giao diện\footnote{Chi tiết được trình bày ở Phần II, Vấn đề 2.} sử dụng một carousel duy nhất đặt ở thanh bên trái hoặc phải, giúp giảm mật độ nội dung trên màn hình chính. Ngoài ra, người dùng có kinh nghiệm có thể chuyển sang \textit{chế độ chỉnh sửa} để kéo thả và sắp xếp lại các thành phần trên trang chủ theo nhu cầu cá nhân.
